\section{Introduction}\label{sec:introduction}

\subsection{Motivation}\label{sec:motivation}

LLM retrieval quality is typically measured by F1 alone, yet token cost is
a first-class production constraint---not a research afterthought. Domain-specific
knowledge has latent structure that RAG discards through chunking, and GraphRAG
re-derives structure from text at significant computational expense. But what if
the structure is already known?

% TODO: Expand motivation with concrete cost examples from production deployments

\subsection{The Three Paradigms}\label{sec:paradigms}

Table~\ref{tab:paradigms} summarizes the three retrieval architectures compared
in this benchmark.

\begin{table}[htbp]
\centering
\caption{Three knowledge retrieval paradigms compared in this benchmark.}
\label{tab:paradigms}
\begin{tabular}{llll}
\toprule
\textbf{System} & \textbf{Knowledge Repr.} & \textbf{Retrieval} & \textbf{Build Cost} \\
\midrule
RAG & Unstructured text chunks & Embedding similarity & Embed all chunks \\
GraphRAG & Dynamically extracted graph & Graph + community search & Full entity extraction \\
CKG & Pre-structured DAG + taxonomy & Direct concept/edge lookup & Zero (CSV-native) \\
\bottomrule
\end{tabular}
\end{table}

\subsection{Falsifiable Claims}\label{sec:claims}

We make the following falsifiable claims, each testable against the benchmark
results presented in Section~\ref{sec:results}:

\begin{enumerate}
  \item CKG achieves higher F1 on T2 (dependency) and T3 (multi-hop path) queries.
  \item CKG F1 does not degrade with hop depth; RAG F1 degrades significantly.
  \item CKG RDS ratio $\geq$ 10x vs RAG across all 25 domains.
  \item GraphRAG hallucinates edges not present in ground truth DAG (HR $>$ 0).
  \item CKG Hallucination Rate $=$ 0 (by construction).
  \item The ``Structure Premium'' hypothesis: RDS advantage correlates with DAG richness ($r > 0.7$).
\end{enumerate}

\subsection{Contributions}\label{sec:contributions}

This paper makes the following contributions:

\begin{enumerate}
  \item The CKG architecture specification (format, DAG constraints, taxonomy schema).
  \item Five novel evaluation metrics (RDS, CUR, Hop-F1, CPCA, RP).
  \item The McCreary Corpus as a formal benchmark dataset (first paper to do so).
  \item An open benchmark: 25 domains $\times$ $\sim$175 queries $\times$ 3 systems ($\sim$13,000 evaluated pairs).
  \item A HuggingFace dataset with one-command reproduction harness.
\end{enumerate}
